
\documentclass{article}
\usepackage{colortbl}
\usepackage{makecell}
\usepackage{multirow}
\usepackage{supertabular}

\begin{document}

\newcounter{utterance}

\centering \large Interaction Transcript for game `cladder', experiment `full\_v1.5\_default', episode 3701 with score{-}Qwen3.5{-}9B{-}sp.
\vspace{24pt}

{ \footnotesize  \setcounter{utterance}{1}
\setlength{\tabcolsep}{0pt}
\begin{supertabular}{c@{$\;$}|p{.15\linewidth}@{}p{.15\linewidth}p{.15\linewidth}p{.15\linewidth}p{.15\linewidth}p{.15\linewidth}}
    \# & \multicolumn{2}{c}{Player} && \multicolumn{2}{c}{Game Master} \\
    \hline

    \theutterance \stepcounter{utterance}  
    & & & \multicolumn{4}{p{0.6\linewidth}}{
        \cellcolor[rgb]{0.9,0.9,0.9}{
            \makecell[{{p{\linewidth}}}]{
                \texttt{\tiny{[P1$\langle$GM]}}
                \texttt{You are an expert in causal inference. The following question is not a typical commonsense query, but rather a meticulously designed question created by a professor specializing in causal inference, intended to assess the students' mastery of the course content.} \\
\texttt{Imagine a self{-}contained, hypothetical world with only the following conditions, and without any unmentioned factors or causal relationships: Education level has a direct effect on skill. Skill has a direct effect on black hair. The overall probability of college degree or higher is 6\%. The probability of high school degree or lower and black hair is 25\%. The probability of college degree or higher and black hair is 3\%. Is the chance of black hair smaller when observing college degree or higher?} \\
\texttt{Start your answer with 'Yes' or 'No', followed by additional reasoning or evidence to support your explanation.} \\
            }
        }
    }
     \\ \\

    \theutterance \stepcounter{utterance}  
    & \multicolumn{4}{p{0.6\linewidth}}{
        \cellcolor[rgb]{0.561,0.737,0.561}{
            \makecell[{{p{\linewidth}}}]{
                \texttt{\tiny{[P1$\rangle$GM]}}
                \texttt{No} \\
\\ 
\texttt{**Reasoning:**} \\
\\ 
\texttt{To determine if the chance of black hair is smaller when observing a college degree or higher, we must compare the conditional probability \$P(\textbackslash{}text\{Black Hair\} \textbackslash{}mid \textbackslash{}text\{College\})\$ with the marginal probability \$P(\textbackslash{}text\{Black Hair\})\$.} \\
\\ 
\texttt{1.  **Calculate the marginal probability of Black Hair, \$P(\textbackslash{}text\{Black Hair\})\$:**} \\
\texttt{    We are given the probability of the intersection of "High School or lower" and "Black Hair" (\$P(\textbackslash{}text\{HS\} \textbackslash{}cap \textbackslash{}text\{BH\}) = 0.25\$) and the probability of the intersection of "College" and "Black Hair" (\$P(\textbackslash{}text\{College\} \textbackslash{}cap \textbackslash{}text\{BH\}) = 0.03\$).} \\
\texttt{    Assuming these are the only two mutually exclusive education categories in this hypothetical world (College or High School/Lower), the total probability of Black Hair is the sum of these intersections:} \\
\texttt{    \$\$P(\textbackslash{}text\{Black Hair\}) = P(\textbackslash{}text\{HS\} \textbackslash{}cap \textbackslash{}text\{BH\}) + P(\textbackslash{}text\{College\} \textbackslash{}cap \textbackslash{}text\{BH\})\$\$} \\
\texttt{    \$\$P(\textbackslash{}text\{Black Hair\}) = 0.25 + 0.03 = 0.28\$\$} \\
\\ 
\texttt{2.  **Calculate the conditional probability of Black Hair given College, \$P(\textbackslash{}text\{Black Hair\} \textbackslash{}mid \textbackslash{}text\{College\})\$:} \\
            }
        }
    }
    & & \\ \\

    \theutterance \stepcounter{utterance}  
    & & & \multicolumn{2}{p{0.3\linewidth}}{
        \cellcolor[rgb]{0.9,0.9,0.9}{
            \makecell[{{p{\linewidth}}}]{
                \texttt{\tiny{[GM$|$GM]}}
                \texttt{No} \\
            }
        }
    }
    & & \\ \\

    \theutterance \stepcounter{utterance}  
    & & & \multicolumn{2}{p{0.3\linewidth}}{
        \cellcolor[rgb]{0.9,0.9,0.9}{
            \makecell[{{p{\linewidth}}}]{
                \texttt{\tiny{[GM$|$GM]}}
                \texttt{no} \\
            }
        }
    }
    & & \\ \\

    \theutterance \stepcounter{utterance}  
    & & & \multicolumn{2}{p{0.3\linewidth}}{
        \cellcolor[rgb]{0.9,0.9,0.9}{
            \makecell[{{p{\linewidth}}}]{
                \texttt{\tiny{[GM$|$GM]}}
                \texttt{game\_result = WIN} \\
            }
        }
    }
    & & \\ \\

\end{supertabular}
}

\end{document}
